\chapter{黎曼量子力学}

在勾股定理符号反转的黎曼宇宙（Riemannian universe）中，物质能否稳定存在？经典图像给出的答案是否定的：处于势能谷（energy valley）中的带电粒子一旦振动就会辐射，而黎曼物理中动能与辐射能的符号相反，辐射反而使它加速，最终逃逸出谷。本章用量子力学来回答这一问题。我们先分析晶体点阵（crystal lattice）中单个粒子所感受到的、由无穷多近邻叠加而成的势井（energy well），再讨论非完美谐振子（imperfect harmonic oscillator）的经典困境、固体中的量子效应与稳定性判据，进而考察两粒子"黎曼原子"（Riemannian atom）与可能的大分子结构；章末以两节进阶材料给出本书采用的单位约定，以及黎曼相对论量子力学中能量与动量算符的构造。

\section{固体的结构（The Structure of Solids）}

在我们的宇宙中，静电库仑力（Coulomb force）服从简单的平方反比律（inverse-square law）：同种电荷相斥、异种电荷相吸，故正电荷附近的正电荷会感受到一个形如陡峭光滑山坡的“能量地形（energy landscape）”，其坡度随距离逐渐变平；负电荷附近则是一个能量坑（energy pit），即把前者上下颠倒。

在黎曼宇宙中，相应的力在近处使\emph{同种电荷相吸、异种电荷相斥}。但这只在某一距离以内成立，因为荷电粒子周围的能量地形之上，叠加了在平方反比衰减上规则排列的一串波峰与波谷。

\orthofig{images/04__003.png}{正电荷周围黎曼与洛伦兹静电势}
\orthofig{images/04__004.png}{负电荷周围黎曼与洛伦兹静电势}

右图给出能量随距离的变化。在中心粒子周围的三维空间中，这些波峰与波谷将形成同心球壳（concentric shells），从而粒子有可能落入彼此的能量波谷而被束缚在一起。

但其中包含一个麻烦。若用经典物理分析（即不计量子力学效应），能量谷中的粒子不会静止，因为它通常拥有足够能量来回运动；而这样运动就会发射电磁辐射。在我们宇宙中，这会使它损失动能、运动衰减；但在黎曼宇宙中，动能与电磁辐射能的符号相反，因此在能量谷中振荡的粒子在辐射时反而会\emph{加速}。若它这样不断获得动能，最终会彻底逃出谷。

粒子得以保持束缚的途径之一是：它振荡的固有频率大于电磁辐射的最大频率 $\nu_{max}$。带电粒子以高于 $\nu_{max}$ 的频率振荡时不会产生辐射。

然而，要产生单一纯频率的振荡，能量谷必须具有完美的形状，即势能作为到中心距离的函数是一条无瑕的抛物线（parabola）。这是非常苛刻的要求，库仑力当然无法满足。对“不完美”的能量谷，即使振荡的主频率高于 $\nu_{max}$，粒子运动也仍含有较低频率的分量。

那么，黎曼宇宙中任何种类的物质如何保持稳定？答案在于量子力学。稍后我们会讨论细节，但为此需要构造一个简单而具体的例子来演示这一效应。

与其考察只有两个粒子由静电力束缚的黎曼版氢原子，不如从另一极端入手，考虑一个粒子被大量规则排列的近邻所包围时所感受到的能量地形。

我们要特别考察的这种排列有多种名称：晶体学家称之为“面心立方系统（face-centred cubic system）”，数学家则称之为“$A_{3}$ 点阵（$A_{3}$ lattice）”。若取一组立方体网格，在每个立方体的每个顶点处以及每个面的中心放一个粒子，就得到 $A_{3}$ 点阵。数学上，取所有满足 $x+y+z$ 为偶整数的点 $(x,y,z)$，便得到一种每边长为 $2$、最近邻距离为 $\sqrt{2}$ 的点阵。当然，我们可以把整个点阵乘上任意因子 $s$ 随意缩放。

\orthofig{images/07__001.png}{$A_{3}$ 或面心立方点阵}
\orthofig{images/07__002.png}{堆叠的橘子与 $A_{3}$ 结构}

左上图中蓝色平面展示了“切开”同一点阵的另一种方式。这些平面被分成等边三角形，顶点处是粒子。这种排列正是球形物体（如橘子或台球）在一个平面内自然堆叠的方式，而相邻平面彼此契合，就像一层层橘子嵌入下一层的凹陷中一样。

$A_{3}$ 点阵中每个粒子有十二个最近邻：同一“橘子层”内六个，上一层三个，下一层三个。由于最近邻距离为 $\sqrt{2}$，其距离平方为 $2$。在我们最初的方案中，每个点阵点具有坐标 $(x,y,z)$ 且 $x+y+z$ 为偶整数，而由
\begin{equation}
(x+y+z)^{2}-(x^{2}+y^{2}+z^{2})=2(xy+yz+xz)
\end{equation}
可知 $(x,y,z)$ 到 $(0,0,0)$ 的距离平方也必为偶整数。由此，$(0,0,0)$ 的近邻距离依次为
\[
\sqrt{2},\ \sqrt{4},\ \sqrt{6},\ \sqrt{8},\ \sqrt{10}\ \dots
\]
点阵中每一点的近邻都呈相同构型。处于逐渐增大距离上的粒子数目，就是把 $2,4,6,8,10\dots$ 写成三个整数平方和的方式数（允许零与负数，且不同的排列顺序分别计数）。前几项计数为
\[
12,\ 6,\ 24,\ 12,\ 24,\ 8,\ 48,\ 6,\ 36,\ 24\ \dots
\]
但这一数列没有简单的公式。

现在假设我们有一团按 $A_{3}$ 点阵排列的正电荷粒子。若这团只有 $4$ 个粒子，它们可以彼此等距，我们还能选择合适的整体尺度使它们都恰好落在彼此最近邻能量波谷的正中心。但对一大团粒子，粒子间距将包含上面列出的各种近邻距离，即所有偶整数的平方根。显然它们\emph{并不}都是某一取值的整数倍，因此不可能让它们各自都恰好位于能量波谷底部。于是整体点阵间距将是一个折中值：它在“让最近邻距离接近库仑势的谷-谷间距 $\lambda_{min}$”与“这一选择对其余更远近邻的影响”之间权衡能量、使总能量最小。

结果是对非常大的团簇，当
\begin{equation}
s=\frac{\sqrt{3}}{2}\lambda_{min}\approx 0.866\,\lambda_{min}
\end{equation}
时团簇势能最小，而最近邻距离为
\begin{equation}
d=\sqrt{2}\,s=\sqrt{\tfrac{3}{2}}\,\lambda_{min}\approx 1.2247\,\lambda_{min}.
\end{equation}
由于这几乎恰好位于 $\lambda_{min}$ 上方四分之一波长处，最近邻对能量的贡献会非常小！尽管如此，包含更多更远近邻的求和结果显示，在这个 $s$ 取值下能量井达到最深。

此时引入精细结构常数（fine structure constant）$\alpha$ 会很有用，它是一个无量纲量，定义为（在时间与空间采用相同单位、即本书通篇所用单位下）
\begin{equation}
\alpha=\frac{e^{2}}{4\pi\hbar\varepsilon_{0}} .
\end{equation}
这里 $e$ 是构成点阵的粒子的电荷，$\varepsilon_{0}$ 是决定静电力强度的介电常数（permittivity constant），$\hbar$ 是普朗克常数除以 $2\pi$。在我们自己的宇宙中 $\alpha\approx 1/137$，但黎曼宇宙中当然没有理由与之相同。

我们再定义
\begin{equation}
m_{\mathrm{photon}}=\hbar\,\wms .
\end{equation}
给定角频率的光，其每个\emph{量子}（quantum）的能量等于该角频率乘以 $\hbar$。因此 $\hbar\wms$ 就是最大可能角频率的光所对应的能量。而正如我们考察黎曼宇宙中的能量与动量时所见，粒子的最大总能量就是其静止质量，所以 $\hbar\wms$ 正是一个黎曼光子的静止质量。

把这两个常数相乘，乘积 $\alpha m_{\mathrm{photon}}$ 等于相距 $\lambda_{min}$ 的两个粒子势能的 $2\pi$ 倍。这使得它成为任何“粒子间距与 $\lambda_{min}$ 挂钩”情形下一个有用的能量标度。

\orthofig{images/07__003.png}{$A_{3}$ 点阵中井深随粒子数的变化}

右图的双对数图给出了一粒子在一系列近似球形的 $A_{3}$ 团簇中心所感受到的能量井的深度。纵轴以 $\alpha m_{\mathrm{photon}}$ 为单位，横轴以 $\lambda_{min}$ 为单位。

作为对比的红线给出同样电荷\emph{均匀}分布于同一区域时井可能达到的最大深度。把黎曼库仑势对实心球积分所得出的精确深度是振荡的，在这一红线与相反的能量（此时井变成峰）之间来回循环。

由于并未采样每个团簇尺寸，能量随团簇增大的振荡被追踪得有些不均匀——但清晰的趋势是井深与团簇尺寸成正比，平均而言约为均匀分布电荷情形的两倍：
\begin{equation}
D_{centre}\approx \frac{8}{3\sqrt{3}\,\pi}\left(\frac{R}{\lambda_{min}}\right)\alpha\,m_{\mathrm{photon}} .
\end{equation}
这不是我们会在自己宇宙的晶体固体中期待的现象：在那里超过某一点后，晶体尺寸对深处电子所感受的条件毫无影响！人们可能以为这种长程影响必定源于“点阵完全由正电荷构成、没有负电荷抵消其累积效应”。但事实上，在点阵中加入负电荷在这一方面几乎不会有差别：它们将位于自己的势能极小点阵上（这些点在正电荷看来像是能量峰），而它们所产生的势能图样实际上会\emph{加深}正电荷的能量井。

\orthofig{images/07__004.png}{能量井形状与完美抛物线的对比}

能量井的深度很重要，但形状呢？任何有限团簇的井都不会完美地旋转对称，且在点阵内各处也不相同；但考虑一个理想化的旋转对称井——势能只依赖于到中心的距离——仍然有用。为近似井的形状，我们可以利用均匀带电球壳内部势能曲线的结果：总势是许多这样一项之和，而每一项都是下式的常数倍
\begin{equation}
U_{s}(r)=-\frac{\sin(\wms r)}{\wms r}.
\end{equation}
既然我们能把点阵近似成一系列带电球壳，中心井的形状必是这一 $U_{s}(r)$ 的某个固定倍数。

$U_{s}(r)$ 本身当然不是抛物线。右图中的蓝线给出 $U_{s}(r)$，而红线是在 $r=0$ 处与它具有相同二阶导数的抛物线。即使把它理想化为旋转对称，井形状的“不完美”也意味着：在这口井中来回运动的经典粒子不会只有单一纯振动频率。因此即使我们保证其振荡主频率大于 $\nu_{max}$，在经典物理中也还不足以阻止它通过运动的较低频率分量向外辐射。

\section{不完美谐振子（Imperfect Harmonic Oscillators）}

当粒子在一口能量井中运动、且（相对井底的）势能正比于到井中心距离的\emph{平方}时，把它拉向中心的力就简单地正比于距离。这样的系统称为谐振子（harmonic oscillator）。弹簧上振动的重物、重力下摆动的单摆都是例子——不过与多数真实物理系统一样，其能量井只是近似抛物，振荡较小时描述最佳。

假设有一口完美抛物、三维径向对称的能量井，粒子在其中无摩擦、无辐射、无任何复杂因素地运动。若力与距离之间的比例常数为 $k$，则势能 $U$ 与力 $\bF$ 为
\begin{align}
U&=\half k\,r^{2},\\
\bF&=-\partial_{r}U\,\be_{r}=-k\,r\,\be_{r}.
\end{align}
把这一径向力分解为沿 $x$、$y$、$z$ 轴的分量。若粒子坐标为 $(x,y,z)$，则径向单位矢量为
\begin{equation}
\be_{r}=\left(\frac{x}{r},\frac{y}{r},\frac{z}{r}\right),
\end{equation}
于是力的各分量为
\begin{equation}
(F_{x},F_{y},F_{z})=(-kx,-ky,-kz).
\end{equation}
利用“力等于质量乘加速度”求每个坐标对时间的二阶变化率，得
\begin{equation}
\bigl(\partial_{t}^{2}x(t),\partial_{t}^{2}y(t),\partial_{t}^{2}z(t)\bigr)
=\bigl(-\tfrac{k}{m}x(t),-\tfrac{k}{m}y(t),-\tfrac{k}{m}z(t)\bigr).
\end{equation}
这三个方程的解为
\begin{equation}
(x(t),y(t),z(t))=
\bigl(A_{x}\cos(\omega t+\phi_{x}),\,
A_{y}\cos(\omega t+\phi_{y}),\,
A_{z}\cos(\omega t+\phi_{z})\bigr),
\end{equation}
其中 $\omega=\sqrt{k/m}$。

\orthofig{images/07__005.png}{抛物与非抛物能量井中的轨道}
\orthofig{images/07__006.png}{做轨道运动的粒子的 $y$ 坐标}

粒子将沿每个坐标方向以相同的频率 $\omega$ 来回振荡。振荡振幅 $A_{x}$、$A_{y}$、$A_{z}$ 一般不同，相位 $\phi_{x}$、$\phi_{y}$、$\phi_{z}$ 也各异。但尽管这些参数有差异，由于\emph{频率}相同，经过时间 $2\pi/\omega$ 后粒子总会精确回到同一点。它走一条\emph{闭合轨道}（closed orbit），一般情形下是一个以井中心为中心的椭圆，如右图中的红色曲线。

若能量井不严格是抛物型呢？力将不再严格正比于 $r$，且力的各分量也不再只依赖于对应坐标。一切都变得更复杂，唯一闭合的轨道是完全没有径向运动的圆轨道。右上图中的蓝色曲线给出能量井为 $U_{s}(r)$（我们用于点阵中粒子势能的函数）时的一段轨道。

右下图为粒子运动的 $y$ 分量随时间的变化，以及描述该运动频率成分的傅里叶变换（Fourier transform）。轨道会进动（precession）——即离井中心最远的点随时间旋转，其角频率记为 $\omega_{p}$——这一事实把频谱中最强的峰从抛物井下的单一值 $\omega_{o}$（轨道频率）分裂为两个峰，位于 $\omega_{o}\pm\omega_{p}$。

现在，这些峰又高又窄（注意这些分量的强度是按对数标度绘制的！），因此若 $\omega_{o}-\omega_{p}>\wms$，粒子的运动就只有极少量处于能够辐射的频率上。但“极少量”并非零，随着时间推移粒子仍会获得能量，最终被驱出井。

\section{固体中的量子效应（Quantum Effects in Solids）}

到目前为止我们的处理完全是经典的，但要真正理解这一情形，需要用量子力学。黎曼宇宙中运动不太快的粒子的经典物理由牛顿力学描述，与在我们的宇宙中一样；而由于\emph{非相对论}量子力学是牛顿力学的量子版本，我们可以直接使用它的许多结果。特别地，对抛物能量井，量子谐振子的行为与我们宇宙中完全相同。

\orthofig{images/07__007.png}{能量本征函数}

量子力学中，粒子的态不是用确定的位置与速度描述，而是用波函数（wave function）。它给粒子可能所在的每个空间点联系一个复数，称为\emph{振幅}（amplitude）；在那里找到粒子的概率正比于振幅的模方。

具有确定能量的波函数，其形状随时间保持不变。右图给出一些例子。完整的波函数在三维空间中变化；这里只画出它们对到井中心距离 $r$ 的依赖关系。

数 $n$ 与 $l$ 是\emph{量子数}（quantum number），用来区分三维量子谐振子的各种可能波函数。量子数 $n$ 与粒子总能量相关：
\begin{equation}
E_{n}=\left(n+\tfrac{3}{2}\right)\hbar\omega .
\end{equation}
这里 $n$ 是非负整数，$\hbar$ 是普朗克常数除以 $2\pi$，$\omega$ 是等效经典谐振子的振荡角频率。于是粒子具有一串等间距的\emph{能级}（energy level）。最低可能的能量态 $n=0$ 并不位于井底，而是位于井底上方 $\tfrac{3}{2}$ 个 $\hbar\omega$ 处，其后所有能级间距均为一个 $\hbar\omega$。

量子数 $l$ 描述粒子具有的总角动量。从波函数图可见，随 $l$ 增大，粒子更可能出现在离井中心更远之处。然而谐振子的特别之处在于：$l$ 对粒子的总能量没有任何影响。量子数 $l$ 被限制在 $n$ 与 $0$ 之间，且当 $n$ 为偶数（奇数）时 $l$ 也为偶数（奇数）。第三个量子数通常记为 $m$，可取从 $-l$ 到 $l$ 的任意整数值，它规定沿某个选定坐标轴的角动量大小；$m$ 对能量与径向波函数都没有影响，但它与 $l$ 一起决定了波函数在偏离该轴不同角度上的变化方式。

总而言之，三维谐振子中粒子可占据 $(n+1)(n+2)/2$ 个不同的量子态，它们全部具有相同的能量 $E_{n}$。

\orthofig{images/07__008.png}{固体点阵井中的能级}

现在，当能量井不严格是抛物型时，所有波函数的形状都会有微小改变。其最重要的后果是：角动量量子数 $l$ 现在\emph{确实}影响粒子的总能量。于是抛物情形下由单一值构成的能级，现在会分裂成若干个彼此靠近的能级。

右图给出若干粒子质量与井深下的能级。每种情形下的红色能级对应完美抛物能量井；蓝色能级对应同样深度、形状为 $U_{s}(r)$（我们预期点阵中粒子势能的近似形状）的井。

量子力学中，带电粒子不再像经典物理描述的那样连续辐射。相反，若它处于能量为 $E_{A}$ 的某个态，而另有一个能量为 $E_{B}$ 的态可供其跃迁，它可以发射一个能量等于两者之差的光子，并从态 $A$ 跳到态 $B$。在我们宇宙中，原子中的电子就是这样发光的。但在黎曼物理中有两点不同。首先，若 $E_{A}$ 与 $E_{B}$ 采用“动能为正”的约定，则光子能量符号相反。因此只有当 $E_{B}>E_{A}$ 时才会发生跃迁，带电粒子在发射光子的同时跳到\emph{更高}的能级。

另一点不同是，黎曼物理中存在最大可能光子能量 $m_{\mathrm{photon}}$ 需要考虑。若 $E_{B}-E_{A}>m_{\mathrm{photon}}$，则单个光子无法携带足够能量促成跃迁。原则上粒子可以发射两个、三个或四个光子，但所需光子数越多，跃迁概率就越小。

结果是：若能量井中粒子所处的量子态，其上下一档态比它高出若干个最大光子能量，它的处境就是\emph{稳定的}。量子力学中很少有情形能永远稳定，但能隙越大，现状就越稳定。例如，图中第一口井里处于最低能量态的粒子，其上方能隙为 $5.2\,m_{\mathrm{photon}}$，因此需要同时发射\emph{六个}光子才能跃迁到更高态。

关键在于：从抛物能量井变成我们在点阵中预期的那种不完美形状，未必会破坏这种稳定性。能级会有一些移动与分裂，但对许多井而言，某些能级上方仍存在若干个最大光子能量的能隙。某个这样的能级\emph{下方}即使存在小得多的能隙也无妨——粒子可以从环境中吸收一个光子下到更低能级，再重新发射光子跳回原处；这类跃迁只是让固体与光发生相互作用，不会对其稳定性产生不利影响。（事实上，即便能隙小于 $m_{\mathrm{photon}}$，也存在一些微妙问题决定哪些跃迁能用单个光子完成。）

究竟哪些固体会在量子力学意义下稳定，取决于诸多参数，包括构成点阵的粒子质量以及点阵中能量井的深度。井深取决于点阵尺寸、黎曼静电力的强度以及每个粒子的电荷。我们此前的计算给出
\begin{equation}
D_{centre}\approx \frac{8}{3\sqrt{3}\,\pi}\left(\frac{R}{\lambda_{min}}\right)\alpha\,m_{\mathrm{photon}},
\end{equation}
其中 $R$ 是团簇半径，$D_{centre}$ 是团簇正中心处井的深度。我们将修改这一公式以计入两点小麻烦。第一，我们希望确保连团簇边缘的井也稳定，为此先把这一估计折半。第二，我们希望允许每口井中有多个粒子存在。因此把估计乘以 $W$，即每口井中的粒子数，因为每个粒子对贯穿整个点阵的场都贡献大致相同的份额
\begin{equation}
D\approx \frac{4W}{3\sqrt{3}\,\pi}\left(\frac{R}{\lambda_{min}}\right)\alpha\,m_{\mathrm{photon}} .
\end{equation}
当然，与我们宇宙中一样，一个现实的固体可能由许多小晶体组成，而非从一端到另一端保持完美点阵结构；此时 $R$ 应是其中一个晶体的半径，而非整块固体的尺度。

设点阵中典型能量井的深度为 $D$。势的最佳抛物近似为
\begin{equation}
U_{p}(r)=D\left[\frac{1}{6}\wms^{2}r^{2}-1\right].
\end{equation}
除了共享深度 $D$ 外，此函数在 $r=0$ 处还与点阵势 $U_{s}(r)$ 具有相同的二阶导数。此井中经典粒子所受的力为
\begin{equation}
\bF=-\partial_{r}U_{p}(r)\,\be_{r}=-\frac{D}{3}\wms^{2}r\,\be_{r},
\end{equation}
于是力常数 $k$ 等于 $\tfrac{D}{3}\wms^{2}$，质量为 $m$ 的粒子在井中的固有振荡频率为
\begin{equation}
\omega=\sqrt{\frac{k}{m}}=\wms\sqrt{\frac{D}{3m}} .
\end{equation}
这意味着井中能量台阶大小 $\hbar\omega$ 与最大光子能量 $m_{\mathrm{photon}}=\hbar\wms$ 之比为
\begin{equation}
\frac{\hbar\omega}{m_{\mathrm{photon}}}=\frac{\omega}{\wms}=\sqrt{\frac{D}{3m}} .
\end{equation}
因此，忽略能级的移动与分裂，点阵的稳定性将取决于 $\sqrt{D/(3m)}$ 的取值。若 $\sqrt{D/(3m)}\le 1$，点阵将\emph{不}稳定；而这个量越大，点阵越稳定。右图说明了这一效应：当粒子质量 $m$ 增大或 $D$ 减小时，不同 $n$ 对应的能级靠得更近。

若我们希望保证
\begin{equation}
\frac{\hbar\omega}{m_{\mathrm{photon}}}=\sqrt{\frac{D}{3m}}>S
\end{equation}
对某个“安全因子（safety factor）”$S$ 成立，则粒子质量与光子质量之比必须小于某个临界值，即 \textbf{(1a)}。若该比值固定，则只要点阵尺寸大于 \textbf{(1b)} 给出的阈值，就能达到同样水平的稳定性。

\begin{keybox}[点阵稳定性条件（Lattice stability conditions）]
\begin{equation}
\frac{m}{m_{\mathrm{photon}}}<\frac{4W\alpha}{9\sqrt{3}\,\pi S^{2}}\left(\frac{R}{\lambda_{min}}\right),\tag{1a}
\end{equation}
\begin{equation}
\frac{R}{\lambda_{min}}>\frac{9\sqrt{3}\,\pi S^{2}}{4W\alpha}\left(\frac{m}{m_{\mathrm{photon}}}\right).\tag{1b}
\end{equation}
\end{keybox}

粒子质量较小有利于稳定性，因为这增大了能级间距，需要更多光子才能跨越该间距。但质量也可能\emph{过}小。完美谐振子的能量井会一直延伸下去，井壁越升越高；而点阵中的井宽度与深度都有限。若能隙相对于井深过大，井中占据最高能级的粒子就有跳出井（而非跃迁到更高的束缚态）的风险。事实上，若能级间距过大，井中甚至可能装不下我们所需的 $W$ 个态。

假设构成点阵的粒子是费米子（fermion），因此与我们宇宙中的电子一样不能共占同一量子态，那么这 $W$ 个粒子将填满若干能级。量子谐振子中，对每个量子数 $n$ 有 $(n+1)(n+2)/2$ 个不同态；若再计入不同自旋的态则数量翻倍，因此到并包括任一给定 $n$ 为止的态总数为
\begin{equation}
\sum_{i=0}^{n}(i+1)(i+2)=\frac{(n+1)(n+2)(n+3)}{3}.
\end{equation}
对总共 $W$ 个粒子所占据的最高能级精确求解非常繁琐，但我们可以作粗略近似，取
\begin{equation}
n\approx \sqrt[3]{3W}.
\end{equation}
“第 $n$ 个能级能装进井内”这一要求给出了质量比的下限，即下面的 \textbf{(1c)}，以及对应的 $R$ 下限 \textbf{(1d)}。

现在，随着点阵增大，质量下限 \textbf{(1c)} 越来越低，而质量上限 \textbf{(1a)} 越来越高；因此无论 $\alpha$ 与 $S$ 取何值，足够大的点阵都会稳定。反过来，对足够小的点阵，我们会到达上限与下限相遇之处，此时两个条件无法同时满足。条件 \textbf{(1e)} 就是要求点阵足够大、以至于不会出现这种情况。

\begin{keybox}[点阵稳定性条件（续）（Lattice stability conditions{,} continued）]
\begin{equation}
\frac{m}{m_{\mathrm{photon}}}>\frac{3^{7/6}\pi}{4\alpha W^{1/3}}\left(\frac{\lambda_{min}}{R}\right),\tag{1c}
\end{equation}
\begin{equation}
\frac{R}{\lambda_{min}}>\frac{3^{7/6}\pi}{4\alpha W^{1/3}}\left(\frac{m_{\mathrm{photon}}}{m}\right),\tag{1d}
\end{equation}
\begin{equation}
\frac{R}{\lambda_{min}}>\frac{3^{11/6}\pi S}{4W^{2/3}\alpha}.\tag{1e}
\end{equation}
\end{keybox}

\orthofig{images/07__009.png}{点阵的稳定性条件}

举个例子，若 $\alpha=1/137$（我们自己宇宙中精细结构常数的值），且 $S=4$（即确保所有粒子都需同时发射至少四个光子才能跳出其基态），则条件 \textbf{(1e)} 要求
\begin{equation}
\frac{R}{\lambda_{min}}>\frac{3225}{W^{2/3}}.
\end{equation}
右图给出稳定点阵可能范围的粗略图景。我们忽略了能级分裂的效应（对最高能级它可能起失稳作用），也忽略了同一井中粒子之间的相互作用（它可能通过降低某些能级而改善稳定性）。尽管如此，此处显现的总体趋势应当成立：对给定粒子质量，\emph{向每口井中加入更多粒子}能使更小的点阵保持稳定。

\section{黎曼原子？（Riemannian Atoms?）}

讨论了含数百万粒子的点阵之后，现在来分析只用两个全同粒子构成的系统。原则上，它们可能以两种不同的方式束缚在一起。它们可能相隔大致 $\lambda_{min}$ 的距离，使各自处于对方的第一个“能量波谷”中；或者它们可能比 $\lambda_{min}/4$ 更近，从而处于彼此的“能量坑”中——即在静电力开始在吸引与排斥之间往复切换之前出现的、无底的静电吸引井。

但从我们对点阵的计算可以看出，若一对粒子处于彼此的能量波谷中，它们不太可能稳定。尽管具体情形与点阵不同，影响能量井深度与能级间距的一般因素却是相同的。对 $W=1$、$R=\lambda_{min}$ 的条件 \textbf{(1e)} 变为
\begin{equation}
\alpha>\frac{3^{11/6}\pi}{4}S\approx 5.89\,S .
\end{equation}
即便假设我们可以自由挑选黎曼宇宙中所有基本常数的值，若为了让此式成立而选取很大的 $\alpha$，那只会增加粒子在不同能级间跃迁的\emph{速率}，从而又要求我们选取更大的 $S$ 以保证稳定。

不过，若两个粒子靠得近得多，我们就能完全忽略库仑力中的波纹，把静电吸引当作平方反比力处理。这类量子力学系统的能级非常著名，因为除若干常数改变外，它们就是氢原子的能级。当然氢原子由一个带负电的电子被束缚在重得多、带正电的质子周围，但由平方反比吸引所决定的波函数数学在两种情形下本质相同。

若我们套用氢原子的量子力学，会发现最简单的两粒子黎曼原子的能级为
\begin{equation}
E_{1}=-\frac{m\alpha^{2}}{4},\qquad
E_{n}=\frac{E_{1}}{n^{2}} .
\end{equation}
这里 $m$ 是每个粒子的质量，$\alpha$ 是精细结构常数，量子数 $n$ 是大于等于 $1$ 的整数。对氢原子原则上 $n$ 没有上限，但对我们的黎曼原子，若 $n$ 过大，波函数将延伸到能量坑之外，我们所用的平方反比力近似就失效。事实上，我们至少需要到 $n=2$ 这一近似依然合理。波函数的距离尺度可通过套用氢原子的玻尔半径（Bohr radius）得到；我们场景下的合适取值为
\begin{equation}
a_{0}=\frac{2\hbar}{m\alpha}.
\end{equation}
要使 $n=2$ 的波函数足够局域化，需要
\begin{equation}
a_{0}<\frac{\lambda_{min}}{64}.
\end{equation}
这里的因子 $1/64$ 有些随意，但它使粒子间距大于 $\lambda_{min}/4$（势首次过零之处）的概率低于 $0.1\%$。这对应一个关于粒子质量与光子质量之比的条件
\begin{equation}
\frac{m}{m_{\mathrm{photon}}}>\frac{64}{\pi\alpha}.\tag{2a}
\end{equation}

那么，这个原子处于 $n=1$ 态时保持稳定的条件是什么？第一与第二能级之差为
\begin{equation}
\Delta E=E_{2}-E_{1}=-\frac{3}{4}E_{1}=\frac{3m\alpha^{2}}{16}.
\end{equation}
关键在于这一能量差包含多少个光子质量：
\begin{equation}
\frac{\Delta E}{m_{\mathrm{photon}}}=\frac{3\alpha^{2}}{16}\left(\frac{m}{m_{\mathrm{photon}}}\right).
\end{equation}
为使原子免于自发发射光子，我们要求此比值超过某个安全因子 $S$，于是
\begin{equation}
\frac{m}{m_{\mathrm{photon}}}>\frac{16S}{3\alpha^{2}}.\tag{2b}
\end{equation}
对 $S=4$、$\alpha=1/137$，这意味着 $m/m_{\mathrm{photon}}>400{,}000$！所以两粒子原子所需的粒子质量，远大于使尺寸合理的点阵稳定的那类粒子。

若原子中不是两个粒子而是 $Z$ 个粒子，能级将乘以 $Z^{2}$。然而，假设这些粒子是费米子，则这 $Z$ 个粒子将填满若干能级。量子数为 $n$ 的不同量子态有 $2n^{2}$ 个，因此最高的 $n$ 值由下式给出
\begin{equation}
Z=2(1+4+9+\dots+n^{2})=\frac{n(n+1)(2n+1)}{3}\approx\frac{2n^{3}}{3},
\end{equation}
\begin{equation}
n\approx\sqrt[3]{\frac{3Z}{2}}.
\end{equation}
（注意这些计算\emph{仅对大 $Z$ 是粗略近似}，不能用于我们的两粒子原子。）我们需要的不是在 $n=1$ 与 $n=2$ 之间留出安全能隙，而是在 $n$ 与 $n+1$ 之间：
\begin{align}
\Delta E&=Z^{2}(E_{n+1}-E_{n})\nonumber\\
&=Z^{2}(-E_{1})\left[\frac{1}{n^{2}}-\frac{1}{(n+1)^{2}}\right]\nonumber\\
&=Z^{2}\frac{m\alpha^{2}}{4}\frac{2n+1}{(n+1)^{2}n^{2}}\nonumber\\
&\approx 2Z^{2}\frac{m\alpha^{2}}{4n^{3}}\nonumber\\
&\approx \frac{Zm\alpha^{2}}{3}.
\end{align}
要求对此能隙有安全因子 $S$，意味着
\begin{equation}
\frac{m}{m_{\mathrm{photon}}}>\frac{3S}{Z\alpha^{2}}.\tag{3}
\end{equation}
对 $S=4$、$\alpha=1/137$，即便 $Z=1000$ 的原子也要求 $m/m_{\mathrm{photon}}>225$。

\section{大分子（Macromolecules）}

在固体晶体点阵与黎曼原子这两种情形之间，应当存在一系列可由任意给定质量粒子构成的结构。稳定性的最明显判据是：单个粒子的能级中包含上方有足够能隙将其隔离的态。我们已经看到，通过种种手段加深能量井可以做到这一点，但这样做可能需要把成百上千个粒子纳入结构之中。

为使粒子系统稳定，系统整体的能级也有必要具有同样的能隙，从而使“分子”的集体振动、弯折与旋转同样免于失控的光发射。转动态的能级正比于相应量子数的平方；由于整数平方之间的间距越来越大，转动能谱最终会在每个态上方形成极大能隙。振动与弯折运动则更为棘手：除了要求分子每个自由度各自的能谱有合适能隙外，还需要结构具有足够对称性，使整份能级清单不至于变成一片看似随机的、彼此紧挨的密集谱线——一组一千个全同谐振子、能隙都恰好为 $\Delta E$，其能谱的最小能隙就是 $\Delta E$；但一组一千个彼此能隙各异的谐振子，则可能通过让一个振子跃迁到更高能级、同时让另一个落到更低能级，从而获得比任何单个振子的量子更小的能量。

要精确刻画哪些分子能对所有这些威胁都保持稳定，将是计算量子化学中的一大壮举，因此我们只能勾画一个大致图景。大型的、类聚合物的、甚至可能呈片状的分子会有最好的机会；其结果可能是：除稀薄气体外，最常见的物质状态都具有相对高度长程有序。于是，黎曼宇宙中相当多的化学过程可能并非如我们熟悉的、在随机相互作用的小分子沸腾液体中进行，而更像在液晶（liquid crystal）中发生。点阵内部可见的那种规则能量井排列，实际上会延伸到点阵之外一段距离，许多稳定大分子也可能如此；它们的相互作用将表现为并排沿着能量地形中的锯齿彼此棘轮式地滑动、相互出入。

有关本主题的额外材料，提供给不介意稍高数学门槛的读者。

\section{进阶：单位（Units）}

本书通篇采用时间与距离用相同单位度量的体系。这等价于在我们宇宙中令光速 $c$ 等于 $1$（例如，以米度量距离，并用光走一米所需的时间作为相应的时间单位）。在黎曼宇宙中，这相当于选取时间单位使得勾股定理（Pythagoras's Theorem）成立，即使三角形的一边牵涉的是时间间隔而非空间。在小说《正交》（\emph{Orthogonal}）中，人们经验性地发现，这等同于令蓝光的速度等于 $1$。

在本节中我们将更进一步，选取质量与能量的单位使“约化普朗克常数” $\hbar=h/(2\pi)$ 等于 $1$。此时质量与能量用具有长度倒数（或空间频率）——也就是时间倒数（或时间频率）——量纲的单位来度量。

我们这一特殊选择意味着，频率 $\nu$ 与能量之间的普朗克关系 $E=h\nu$ 变为 $E=2\pi\nu=\omega$，其中 $\omega$ 是波的角频率；空间频率 $\kappa$ 与动量之间的关系则变为 $p=2\pi\kappa=k$，其中 $k$ 是角空间频率。黎曼标量波方程（Riemannian Scalar Wave equation）中出现的最大角频率 $\wms$ 于是就直接等于相应粒子的静止质量。

\section{进阶：相对论能量与动量算符（Relativistic Energy and Momentum Operators）}

在迄今讨论的非相对论量子力学中，我们只是把通常的薛定谔方程（Schrödinger equation）应用于带电粒子间作用力所对应的势能，其依据是：只要把动能视为正、并选取与之相容的势能符号，黎曼宇宙中的非相对论经典动力学就与牛顿力学完全相同。

对于黎曼\emph{相对论}量子力学，我们需要略作不同的处理。量子力学的通常表述与相应经典力学的哈密顿形式（Hamiltonian form）紧密相关；而在黎曼情形中，哈密顿意义下与各坐标共轭的动量，是同一方向上相对论动量的\emph{相反数}。

在非相对论量子力学中，我们把能量等同于算符 $\ii\partial_{t}$，但把动量的空间分量等同于算符 $-\ii\partial_{x}$、$-\ii\partial_{y}$、$-\ii\partial_{z}$。洛伦兹相对论量子力学中也是如此；但在黎曼情形中这说不通：我们需要把空间与时间坐标等同对待。为了保持与哈密顿经典力学的通常对应，我们将保留对应关系
\begin{equation}
E\to \ii\partial_{t},
\end{equation}
并把相对论动量等同于其它坐标上同类算符，而\emph{不带}通常的负号：
\begin{equation}
p_{x}\to \ii\partial_{x},\qquad p_{y}\to \ii\partial_{y},\qquad p_{z}\to \ii\partial_{z}.
\end{equation}
这也会影响我们用于角动量的算符。若用我们的黎曼相对论动量 $\bp$ 定义三维矢量算符 $\mathbf{L}=\mathbf{r}\times\bp$，它变为
\begin{align}
\mathbf{L}&=(x,y,z)\times(\ii\partial_{x},\ii\partial_{y},\ii\partial_{z})\nonumber\\
&=\ii\bigl(y\partial_{z}-z\partial_{y},\ z\partial_{x}-x\partial_{z},\ x\partial_{y}-y\partial_{x}\bigr).
\end{align}
与线动量一样，它也与通常算符相差一个负号。因此，其各分量之间的对易子（commutator）与非相对论量子力学中的不同：
\begin{equation}
[L_{i},L_{j}]=-\ii\varepsilon_{ijk}L_{k},
\end{equation}
其中列维-奇维塔符号（Levi-Civita symbol）$\varepsilon_{ijk}$ 当 $(i,j,k)$ 是 $(x,y,z)$ 的偶排列时等于 $1$，奇排列时等于 $-1$，含重复坐标时等于 $0$。同样，此式右端与通常情形相差一个符号。

现在，为使粒子的\emph{内禀}（intrinsic）角动量与这整套方案相容，相应的算符必须与轨道角动量 $\mathbf{L}$ 具有相同的对易代数。这意味着，对自旋 $\tfrac{1}{2}$ 粒子角动量各分量而言，合适的算符不再是泡利自旋矩阵（Pauli spin matrices），而是其\emph{相反数}。细节将随出现而展开；此处只是提醒：要对通常约定中的这一细小扭转有心理准备。
